% Ad Atticum 14.5 (ad-atticum-14-05) --- English translation
% Translated by Claude (Anthropic) from the Latin (Perseus).
% Style guide: STYLE.md.

\ciceroLetterOpener{CICERO ATTICO S. D.}

\ciceroSection{1} I hope that by now things are as we would wish, since you went without food when you were a little unwell; still, I should like to know how you are. Pretty signs, that Calvena takes it badly that Brutus suspects him; nothing pretty about the signs, if it is with the standards that the legions are coming from Gaul. What of those that were in Spain --- what do you suppose? Will they not demand the same? And what of those Annius shipped over? I meant Gaius Asinius --- but a \textit{[Greek: mnemonikon hamartema]}, a slip of memory. From the bath-keeper a \textit{[Greek: phurmos polus]}, a great muddle. As for that conspiracy of Caesar's freedmen, it could easily be put down, if Antony had any sense.

\ciceroSection{2} What a foolish scruple of mine! I refused a legation before the recess of business, for fear of seeming to desert this swollen state of things --- when, if I could in fact apply a remedy to it, I ought not to fail it. But you see the magistrates, if magistrates is what they are; you see, all the same, the tyrant's hangers-on holding commands; you see his armies still in being; you see the veterans at our flank, all things \textit{[Greek: euripista]}, easily blown about; and the men who ought to be not only ringed but made great by the world's protection, only praised and only loved --- and shut indoors within four walls. And so, however things stand for them, they are happy; the state is wretched.

\ciceroSection{3} But I should like to know about Octavius's arrival --- whether there is any rallying to him, any suspicion of \textit{[Greek: neoterismou]}, of a new agitation. I do not think so, but whatever the truth is, I want to know it. I have written this to you as I set out from Astura, on the third day before the Ides.
